Fix integers $k\ge 2$ and $1\le m\le k$.
Define
\[
f_{k,m}(x):=\bigl|\{n\in\mathbb Z_{\ge 0}: n\le x\text{ and }n=a_1^k+\cdots+a_m^k\text{ for some }a_i\in\mathbb Z_{\ge 0}\}\bigr|.
\]
So $f_{k,m}(x)$ counts \emph{distinct} integers $\le x$ representable as a sum of $m$ nonnegative $k$th powers.

Questions (as written):

(Q1) Is it true that for every $\epsilon>0$,
\[
f_{k,k}(x)\gg_\epsilon x^{1-\epsilon}?\]

(Q2) Is it true that if $m<k$, then
\[
f_{k,m}(x)\gg x^{m/k}\quad\text{for all sufficiently large }x?\]

Edge cases.
$f_{k,m}(x)$ is nondecreasing in $x$, and $f_{k,m}(x)\ge 1$ for all $x\ge 0$ since $0=0^k+\cdots+0^k$.
